% Template for reports/latex/model_fidelity.py, filled by tools/latex_report.py.
% Every number below comes from the script, not from this file.
\documentclass{cadstudio}
\usepackage{xcolor}
\usepackage{subcaption}
\title{\VAR{title}}
\subtitle{\VAR{subtitle}}
\documentid{\VAR{document_id}}
\revision{\VAR{revision}}
\date{\VAR{date}}
\newcommand{\swatch}[1]{{\setlength{\fboxsep}{0pt}\fbox{\textcolor[HTML]{#1}{\rule{16pt}{7pt}}}}}
\newcommand{\code}[1]{\texttt{#1}}

\begin{document}
\maketitle

\begin{abstract}
  The CAD models were compared with reference renders of engines produced by other text-to-CAD
  work, and five shortcomings were identified: no surface finishes, colour used as a function
  code, primitives standing in for components, no fasteners, and sharp edges throughout. The shape
  library was extended to remove them and the \VAR{pilot_title} engine was rebuilt as the pilot.
  It now has \VAR{pilot.after.parts} parts in \VAR{pilot.after.groups} sub-assemblies where it had
  \VAR{pilot.before.parts} parts in a flat list, \VAR{pilot.after.fasteners} fasteners where it had
  none, \VAR{pilot.after.rounded} profiles with rounded or chamfered corners, and
  \VAR{pilot.after.finishes} named finishes. No two parts intersect, and every part rebuilt in
  FreeCAD from the same recipes matches the cadgen part.
  \BLOCK{if pending}
  \VAR{upgraded|length} of \VAR{rows|length} models have been upgraded at this revision.
  \BLOCK{else}
  All \VAR{rows|length} models have been upgraded.
  \BLOCK{endif}
\end{abstract}

\section{Introduction}

Four reference images were studied: a Raptor engine shown in the CAD Viewer with material
shading, an exploded radial aircraft engine, and two views of a Raptor in a custom web viewer.
What separates them from the models in this project is not the accuracy of the main contours,
which here follow published dimensions, but everything around those contours. The reference
models have metal that looks like metal, a few real finishes in place of a colour per function,
bolted flanges, rounded edges that catch a highlight, and an assembly tree that opens into
sub-assemblies with many parts.

Five additions answer this. They are described in section~\ref{sec:features} and their effect on
the pilot model is measured in section~\ref{sec:results} against the model as it stood at commit
\code{\VAR{baseline}}, rebuilt from that commit for the comparison.

\BLOCK{if images}
\begin{figure}[!htb]
  \begin{subfigure}{0.49\textwidth}
    \includegraphics[width=\textwidth]{\VAR{images.before}}
    \caption{Before: \VAR{pilot.before.parts} parts}
  \end{subfigure}\hfill
  \begin{subfigure}{0.49\textwidth}
    \includegraphics[width=\textwidth]{\VAR{images.after}}
    \caption{After: \VAR{pilot.after.parts} parts}
  \end{subfigure}
  \caption{\VAR{pilot_title} engine in the viewer's default display, from the same camera
    position, before and after the upgrade.}
  \label{fig:beforeafter}
\end{figure}
\BLOCK{endif}

\section{What was added}
\label{sec:features}

\subsection{Finishes}

Each part previously carried one flat colour and no surface properties, so the viewer's Render
display and the GLB export showed every part as matt plastic. Parts are now assigned a named
finish from one palette, \code{models/lib/materials.py}, which sets base colour, roughness and
metalness (table~\ref{tab:finishes}). The model declares the assignment once and cadgen writes it
to the STEP sidecar and into the GLB materials.

Two rules came out of the first renders. Base colours are kept mid to light, because the dark
greys used before hid the form of the nozzle. Metalness is kept moderate, because with values
near one the default Render lighting turned large surfaces almost black.

The finishes are display properties. They do not assign an alloy, a density or a mass to a part.

\begin{table}[!htb]
  \caption{Finishes used by the \VAR{pilot_title} model.}
  \label{tab:finishes}
  \begin{tabular}{@{}llcS[table-format=1.2]S[table-format=1.2]S[table-format=2.0]@{}}
    \toprule
    Finish & \multicolumn{2}{l}{Base colour} & {Roughness} & {Metalness} & {Parts} \\
    \midrule
    \BLOCK{for row in finishes}
    \VAR{row.name} & \code{\#\VAR{row.hex}} & \swatch{\VAR{row.hex}} & \VAR{"%.2f"|format(row.roughness)} &
      \VAR{"%.2f"|format(row.metalness)} & \VAR{row.parts} \\
    \BLOCK{endfor}
    \bottomrule
  \end{tabular}
\end{table}

\subsection{Rounded and chamfered edges}

Every revolved profile was a polygon with sharp corners. The new \code{lathe} operation revolves
the same polygon with each corner replaced by a tangent arc (figure~\ref{fig:corners}). The radius
is reduced where an adjacent side is short, and corners on the axis and the small bends along a
curved contour are left alone. The operation also takes an arbitrary axis, so valves, duct flanges
and the gas generator are turned parts on their own centre lines in place of plain cylinders.

A small radius on a large diameter is expensive to mesh: rounding the hatbands and nozzle hoops
took the first export of the pilot to about six million triangles. Those large thin rings are
chamfered instead, which gives the same highlight from a conical face.

\begin{figure}[!htb]
  \includegraphics{corners.pdf}
  \caption{Profile of the \VAR{pilot_title} turbopump bearing housing (\code{\VAR{profile_name}}),
    as written in the recipe and as revolved. Axial position $a$ and radius $r$ are about the
    turbopump axis.}
  \label{fig:corners}
\end{figure}

\subsection{Fasteners}

No model had a fastener. Two pattern operations, \code{hex\_circle} and \code{pin\_circle}, place
hexagonal prisms and round pins on a circle about any axis, and \code{flange\_bolts} combines them
into a bolt circle: a nut and a stud end standing on a flange face. Nothing passes through the
flange, so no holes are cut and the fasteners touch the flange without intersecting it. Each bolt
circle is one part, which keeps the assembly tree and the files small.

\subsection{Sub-assemblies and component detail}

Parts were fused into one solid per function and listed flat. A model can now name
sub-assemblies, and the pilot uses them to split its former single solids into the pieces they
are built from (table~\ref{tab:groups}). The turbopump, formerly one revolved solid, is an inlet,
two volutes, a bearing housing and a turbine with bolted joints between them. The gas generator,
formerly a cylinder, has an injector head, flange, combustor, outlet neck and a duct to the
turbine. The four valves are flanged bodies with actuators.

\begin{table}[!htb]
  \caption{Sub-assemblies of the \VAR{pilot_title} model.}
  \label{tab:groups}
  \begin{tabular}{@{}lS[table-format=2.0]S[table-format=3.0]S[table-format=2.0]S[table-format=1.0]@{}}
    \toprule
    Sub-assembly & {Parts} & {Fasteners} & {Rounded profiles} & {Finishes} \\
    \midrule
    \BLOCK{for row in groups}
    \code{\VAR{row.name}} & \VAR{row.parts} & \VAR{row.fasteners} & \VAR{row.rounded} & \VAR{row.finishes} \\
    \BLOCK{endfor}
    \bottomrule
  \end{tabular}
  \tablenote{\VAR{loose|length} further parts stand at the top level:
    \BLOCK{for label in loose}\code{\VAR{label}}\VAR{", " if not loop.last else "."}\BLOCK{endfor}
    Fasteners are counted as nuts; each has a stud end.}
\end{table}

\subsection{Heat tint}

The tube bundle below the chamber jacket is cut into axial bands, each a separate part with its own
finish, from \qty{\VAR{"%.0f"|format(tint_from)}}{\milli\metre} to
\qty{\VAR{"%.0f"|format(tint_to)}}{\milli\metre} from the gimbal centre. The bands run
\VAR{tints|join(", ")} away from the throat, the order in which oxide colours appear as
temperature falls. This is an appearance choice taken from the reference images. It is not based
on a photograph or a temperature of the F-1.

\section{Results on the pilot model}
\label{sec:results}

Table~\ref{tab:pilot} compares the pilot with its previous state.
\BLOCK{if images}
Figures~\ref{fig:render} and~\ref{fig:detail} show it in the Render display.
\BLOCK{endif}
The number of faces rose by a factor of \VAR{"%.0f"|format(face_ratio)} and the STEP file by a
factor of \VAR{"%.0f"|format(step_ratio)}. The fasteners account for
\VAR{"%.0f"|format(fastener_face_share * 100)}\,\% of the faces; each nut and stud is written as its own solid.

\begin{table}[!htb]
  \caption{\VAR{pilot_title} model before and after the upgrade.}
  \label{tab:pilot}
  \begin{tabular}{@{}lS[table-format=5.1]S[table-format=5.1]@{}}
    \toprule
    Quantity & {Before} & {After} \\
    \midrule
    Parts & \VAR{pilot.before.parts} & \VAR{pilot.after.parts} \\
    Sub-assemblies & \VAR{pilot.before.groups} & \VAR{pilot.after.groups} \\
    Recipe operations & \VAR{pilot.before.operations} & \VAR{pilot.after.operations} \\
    Faces & \VAR{pilot.before.faces} & \VAR{pilot.after.faces} \\
    Fasteners & \VAR{pilot.before.fasteners} & \VAR{pilot.after.fasteners} \\
    Rounded or chamfered profiles & \VAR{pilot.before.rounded} & \VAR{pilot.after.rounded} \\
    Finishes & \VAR{pilot.before.finishes} & \VAR{pilot.after.finishes} \\
    STEP file, MB & \VAR{"%.1f"|format(pilot.before.step_mb)} & \VAR{"%.1f"|format(pilot.after.step_mb)} \\
    GLB file, MB & \VAR{"%.1f"|format(pilot.before.glb_mb)} & \VAR{"%.1f"|format(pilot.after.glb_mb)} \\
    \bottomrule
  \end{tabular}
  \tablenote{The GLB is meshed at a relative chord tolerance of \num{\VAR{mesh_tolerance}} and an
    angular tolerance of \qty{\VAR{mesh_angle}}{\radian}.}
\end{table}

\BLOCK{if images}
\begin{figure}[!htb]
  \includegraphics[width=0.92\textwidth]{\VAR{images.render}}
  \caption{\VAR{pilot_title} model in the Render display with the assigned finishes.}
  \label{fig:render}
\end{figure}

\begin{figure}[!htb]
  \includegraphics[width=0.92\textwidth]{\VAR{images.detail}}
  \caption{Powerhead of the \VAR{pilot_title} model: turbopump housings with bolted joints, duct
    and valve flanges, the heat-tint bands below the chamber jacket and the injector behind
    the oxidizer dome.}
  \label{fig:detail}
\end{figure}
\BLOCK{endif}

\section{Verification}

The saved STEP of each upgraded model was checked with \code{tools/model\_check.py}: every part
must be a valid solid of positive volume, and every pair of parts whose bounding boxes overlap is
intersected with a fuzzy boolean and must share no more than
\qty{\VAR{"%.0f"|format(pilot.after.check.tolerance_mm3)}}{\milli\metre\cubed}.
Table~\ref{tab:checks} gives the results. Where a FreeCAD comparison is listed, the model was
rebuilt in FreeCAD from the same recipes and each part compared with the cadgen part by volume and
overlap.

\begin{table}[!htb]
  \caption{Checks on the upgraded models.}
  \label{tab:checks}
  \begin{tabular}{@{}lS[table-format=3.0]S[table-format=1.0]S[table-format=4.0]S[table-format=1.0]l@{}}
    \toprule
    Model & {Parts} & {Invalid} & {Pairs tested} & {Intersecting} & FreeCAD match \\
    \midrule
    \BLOCK{for row in upgraded}
    \VAR{row.title} & \VAR{row.after.check.parts} & \VAR{row.after.check.invalid|length} &
      \VAR{row.after.check.pairs_tested} & \VAR{row.after.check.intersecting|length} &
      \BLOCK{if row.after.freecad}\VAR{row.after.freecad.ok} of \VAR{row.after.freecad.ok + row.after.freecad.bad} parts\BLOCK{else}not run\BLOCK{endif} \\
    \BLOCK{endfor}
    \bottomrule
  \end{tabular}
\end{table}

Two faults were found and corrected this way during the pilot. A fuel duct flange cut into the
oxidizer volute and was moved. A fused strut returned an inside-out solid without an error, which
removed the chamber jacket from the export; the FreeCAD comparison exposed it, the strut was
reshaped, and \code{build()} now stops when adding a feature reduces the volume of a part.

\section{Status of the models}

Table~\ref{tab:models} lists all models against their state at commit \code{\VAR{baseline}}.
\BLOCK{if pending}
\VAR{pending|length} models have not yet been upgraded and are unchanged.
\BLOCK{endif}

\begin{table}[!htb]
  \caption{All models, before $\rightarrow$ after.}
  \label{tab:models}
  \small
  \begin{tabular}{@{}lccccc@{}}
    \toprule
    Model & Parts & Fasteners & Rounded profiles & Finishes & STEP, MB \\
    \midrule
    \BLOCK{for row in rows}
    \VAR{row.title} &
      \VAR{row.before.parts} $\rightarrow$ \VAR{row.after.parts} &
      \VAR{row.before.fasteners} $\rightarrow$ \VAR{row.after.fasteners} &
      \VAR{row.before.rounded} $\rightarrow$ \VAR{row.after.rounded} &
      \VAR{row.before.finishes} $\rightarrow$ \VAR{row.after.finishes} &
      \VAR{"%.1f"|format(row.before.step_mb)} $\rightarrow$ \VAR{"%.1f"|format(row.after.step_mb)} \\
    \BLOCK{endfor}
    \bottomrule
  \end{tabular}
\end{table}

\BLOCK{if images and images.gallery}
\BLOCK{for item in images.gallery}
\begin{figure}[!htb]
  \includegraphics[width=0.9\textwidth]{\VAR{item.image}}
  \caption{\VAR{item.title} in the Render display.}
\end{figure}
\BLOCK{endfor}
\BLOCK{endif}

\section{Limits}

\begin{itemize}
  \item The added detail is for appearance. Flange sizes, bolt counts, corner radii and the split
    of components into housings are assumed, not taken from the sources, and each model source
    says so. The main contours and envelope dimensions are unchanged.
  \item Finishes appear in the viewer's Render display and in GLB files. The default display
    shows base colours with plain shading.
  \item The models lie along the $x$ axis through the origin, so the Render display's floor at
    $z = 0$ cuts through them. The renders here place the floor at the lowest point of the model.
  \item The new operations build in cadgen and FreeCAD. The SolidWorks builder does not have them.
\end{itemize}

\end{document}
